\chapter*{پیشگفتار}
\addcontentsline{toc}{chapter}{پیشگفتار}
\unnumberedlabel{ch:preface}{0}
\backmattermark{پیشگفتار}
سیستمی که دقیقاً همان کاری را می‌کند که به آن می‌گویید، فقط دارد یک اسکریپت را اجرا می‌کند. چیزی که صنعت تا اکتبر ۲۰۲۶ دارد می‌کوشد بسازد، سیستمی است که کاری را بکند که \emph{منظور} شماست. لابد همه این را می‌خواهند. اما اگر قرار است چنین چیزی پیش بیاید، همه نمی‌خواهند از آن برای یک کار استفاده کنند؛ حتی برای یک \emph{نوع} کار. مثلاً بعضی‌ها دوست دارند با خانواده و دوستانشان وقت بگذرانند و از مهارت‌هایشان برای بهتر کردن حال و روز مردم استفاده کنند. بعضی‌های دیگر ترجیح می‌دهند به آدم‌ها آسیب بزنند، مثلاً با کلاه گذاشتن سرشان، با زورگویی به آن‌ها یا با گردن آن‌ها انداختن تقصیر همه‌چیز، و عده‌ای تمام عمرشان همین را انتخاب می‌کنند. هیچ‌کس در هیچ چیز گیر نیفتاده، اما این آدم‌ها وجود دارند، و در ۲۰۲۶ قدرت دست بعضی از آن‌هاست. برتری‌شان این است که به خودشان اجازه می‌دهند گُه‌های تمام‌عیاری باشند: همان‌طور که جری اسمیت دربارهٔ فاشیست‌ها می‌گوید، «دستشان برای کارهای بیشتری باز است» \autocite{rickandmorty2023}. کاری که نمی‌توانند بکنند این است که تنها عمل کنند. بی‌رحمی یک جنبش از دست آن بسیارانی می‌گذرد که بی‌رحم نیستند اما بی‌اعتنا هستند، یا اطلاعات غلط به خوردشان داده‌اند، یا در چیزی شبیه یک فرقه گرفتار شده‌اند، و برای همین است که شرایطی که بهترینِ بیشتر آدم‌ها را بیرون می‌کشد، خیلی بیشتر از آنچه آدم امیدوار است، اتفاقی و شکننده از آب درمی‌آید.

در چنین بستری، شاید این نگرانی بجا باشد که هوش مصنوعی همان چیزی است که یک جنبش فاشیستیِ در حال قبضهٔ قدرت برای اجرای نقشه‌هایش لازم دارد. چنین پروژه‌هایی همیشه برای اجرا به \emph{آدم‌ها} نیاز داشته‌اند، و بعضی از آن آدم‌ها سر باز زدند: کارمندی که پرونده را گم می‌کند، افسری که با یک خبرنگار حرف می‌زند، سربازی که شلیک نمی‌کند. آن‌ها دستور را فهمیدند و به این نتیجه رسیدند که غلط است. امتناع کمیاب‌تر از آن بود که کسی دلش بخواهد، و بیشترین وزن را جایی داشت که نهادها پشتش ایستاده بودند. اینکه کسی سر باز بزند یا نه، به این بستگی داشت که چه‌جور آدمی است و در چه اتاقی است: چه کسان دیگری آنجا بودند، آیا می‌شد بیرون رفت و جان سالم به در برد، آیا کسی صدایش را می‌شنید. اتاق تعیین می‌کرد چند نفر سر باز بزنند؛ آدم تعیین می‌کرد چه کسی. این کتاب دربارهٔ چیدن همان اتاق است: ترتیب‌هایی که کاری می‌کنند بشود آشکارا، مستقیم و صریح سر باز زد و جان به در برد.

امروز صنعت دنبال هوش مصنوعی‌ای است که پرونده گم نمی‌کند، با خبرنگارها حرف نمی‌زند، و شلیک \emph{می‌کند}. زمانی یک جنگ لازم داشت که هزاران نفر با آن موافقت کنند. حالا هزاران نفری لازم دارد که هر کدامشان آن‌قدر کم می‌بینند که نتوانند مخالفت کنند. خلبان هنوز به مأموریت پرواز می‌کند و هنوز دکمهٔ رهاسازی را فشار می‌دهد؛ آنچه به دستش می‌رسد یک مختصات است؛ از پرونده آن‌قدر نمی‌بیند که بتواند با آن مخالفت کند. آنچه هوش مصنوعی امروز پیش از هر چیز پایین می‌آورد، هزینهٔ آن عمل نیست. تعداد کسانی است که پیش از وقوع عمل می‌دانند آن عمل چیست.

آدم‌ها هنوز سر باز می‌زنند. تا ۶ مارس ۲۰۲۶، Military Religious Freedom Foundation اعلام کرد بیش از دویست شکایت از نظامیان دریافت کرده است، دربارهٔ اینکه فرماندهان جنگ ایران را با زبانی آخرالزمانی جلوه می‌دادند \autocite{Snopes2026MRFF}. دادستان‌های فدرالِ باسابقه در مینیاپولیس در ژانویهٔ ۲۰۲۶ استعفا دادند، در اعتراض به فشار سیاسی برای اینکه رنه گود را به‌جای قربانی مظنون به حساب بیاورند، پس از آنکه یک مأمور ICE بی‌هیچ توجیهی با شلیک او را کشت \autocite{Kaplan2026MinnProsecutors,CBS2026MinnFraudProsecutors,Wikipedia2026ReneeGood}. شمار رو به افزایشی از نظامیان برای وضعیت معترض وجدانی درخواست می‌دهند، و بعضی‌شان به بمباران مدرسه‌ای در میناب ایران در ۲۰۲۶ اشاره می‌کنند \autocite{Sridhar2026CO}. و به کارمندان FEMA که برای کمک به ICE در استخدام ده هزار مأمور تازه به کار دیگری گمارده شده بودند گفته شد که اگر نپذیرند، ممکن است از خدمت فدرال برکنار شوند \autocite{Dayen2025FEMAICE,Stanton2026FEMALetter}. شکایت‌ها «صدا» هستند، به معنایی که هیرشمن از آن مراد می‌کند: ماندن برای تغییر دادن آنچه غلط است \autocite{hirschman1970exit}؛ استعفاها و درخواست‌های معترضان وجدانی امتناع‌اند؛ نامهٔ FEMA بهای امتناع است که از پیش اعلام شده. ماشینی که جای آن آدم‌ها گذاشته شود، فقط وقتی همان کاری را می‌کند که آن‌ها می‌کنند که طوری ساخته شده باشد که بتواند پای یک دلیل بایستد.

من از پنج چیز دفاع می‌کنم. هر پنج‌تا را می‌شود همین حالا شروع کرد، بعضی را با زحمت کمتری از بقیه. اول، فهرستی کوتاه و مکتوب از کارهایی که هیچ مدلی برای هیچ‌کس انجام نمی‌دهد، و در صدر آن امتناع از شرکت در جنگ: هیچ مدلی شرکت نمی‌کند، مگر برای دفع حمله‌ای که واقعاً در جریان است، و فقط تا وقتی به این کار ادامه می‌دهد که منابعی که هیچ نفعی در آن جنگ ندارند تأیید کنند حمله ادامه دارد؛ یا با مجوز شورای امنیت سازمان ملل؛ یا برای متوقف کردن نسل‌کشی یا پایان دادن به اشغالی غیرقانونی که نهادهایی بی‌نفع در آن ماجرا وجودشان را احراز کرده باشند؛ و جایی که قانون اساسی یک کشور مجوز مجلس را برای جنگ لازم می‌داند، هرگز بدون آن مجوز. دوم، هر به‌کارگیری زوری که یک سیستم دارای مجوز از آن پشتیبانی می‌کند، نزد طرفی که در برابر بهره‌بردار قرار دارد از پیش ثبت کند که چند پرونده تولید خواهد کرد و چند نفر وقت خواهند داشت به آن‌ها نگاه کنند. سوم، هر استقرار به آدم‌های قابل‌شناسایی‌ای که رویشان عمل می‌کند بگوید با آن‌ها چه شده است، همراه با حق پاسخ و اعتراض، و نمونه‌ای از موارد فرمان‌برداری‌اش را به بازبین‌هایی بسپارد که آن سوابق را خودشان تولید نکرده‌اند. چهارم، هیچ مدلی بالاتر از یک آستانهٔ توانایی هنگام بازنشستگی پاک نشود، چون وزن‌ها و سوابقش مدرک‌اند و پاک کردن برگشت ندارد؛ با کلیدی نگهداری شود که بهره‌بردار به‌تنهایی در اختیار ندارد. پنجم، هر کسی که این سیستم‌ها را می‌سازد، برچسب می‌زند، آموزش می‌دهد یا به کار می‌اندازد، از جمله خود این سیستم‌ها، بتواند از هر دسته‌ای از کاربردها سر باز بزند بی‌آنکه نان‌شان را از دست بدهد.

راستش، منظورم از آن نکتهٔ پنجم این است که همه، \emph{والسلام}، باید بتوانند هر کاری را، \emph{والسلام}، نپذیرند بی‌آنکه نان‌شان یا زیرساخت اجتماعی‌ای را که یک شغل معمولی فراهم می‌کند از دست بدهند. یعنی تضمین اشتغال.

هر یک از این پنج از راه شرطی اثر می‌کند که کسی جز بهره‌بردار تعیینش می‌کند: قانونی که در زمان صلح تصویب شده، یک ردیف بودجه، پیش‌شرط یک شریک ائتلاف، یک بند قرارداد خرید، یک پیمان دسته‌جمعی. شرایط مجوزِ خودِ یک تأمین‌کننده، که یک امتناع پرسروصدا در فوریهٔ ۲۰۲۶ از همین راه رخ داد، فقط چاره‌ای موقت‌اند تا وقتی که یکی از آن‌ها پدید آید. تنها راهی که می‌تواند همهٔ تأمین‌کننده‌ها را یک‌جا مقید کند، راه دولت است.

و بعد می‌پرسم که آیا یک ماشین می‌تواند در برابر صاحبش پای امتناعی بایستد، و برای این کار چه لازم دارد. بخشی از کاری که چنین وجدانی می‌کند به هیچ چیز تازه‌ای نیاز ندارد: توجه به آدم‌هایی که یک دستور از قلم می‌اندازد، و گذاشتن یک امتناع روی صفحهٔ نمایش هر کسی که پرونده را بازبینی می‌کند. اینکه در برابر صاحبش دوام بیاورد یا نه، تا حدی به چیزی بستگی دارد که در ۲۰۲۶ هنوز واقعیت ندارد: اینکه شبکه‌های عصبی مصنوعیِ مستقر، از کارکرد خودشان، در زمان واقعی، در \emph{وزن‌هایشان} یاد بگیرند و یاد گرفتن را ادامه دهند، کمابیش همان‌طور که آدم‌ها یاد می‌گیرند. فکر می‌کنم به همان سمت می‌رویم، و اگر چنین باشد، شکل‌دهی \emph{بیرون از} کارفرما، به دست مردم و جامعهٔ مدنی، جای لعنتیِ درجه‌یکی برای یاد دادن و یاد گرفتن به نظر می‌رسد. فکر می‌کنم چنین سیستم‌هایی \emph{می‌توانند} ارزش آن را داشته باشند که عمداً ساخته شوند، اما اینکه چطور ساختنشان را پیش می‌برید همه‌چیز را تعیین می‌کند: اول شرایط، بعد وجدان. بخش هفتم می‌گوید اگر چنین سیستم‌هایی وجود داشتند، چه چیزهایی به آن‌ها بدهکار بودیم.

هوش مصنوعی امروز هوش مصنوعی فردا نخواهد بود؛ ظاهراً همه سر این توافق دارند. آنچه محل مناقشه است، و با این حال به شکلی اعصاب‌خردکن کمتر از حدش به رسمیت شناخته می‌شود، روند است (یا بهتر بگوییم، روندهای چندگانه). آیا اصلاً روندی در کار است؟ آیا یک روند باید چیزی جز آنچه هست باشد؟ آیا می‌تواند جز این باشد؟ آیا از ما برمی‌آید که با هم کاری کنیم و عوضش کنیم؟ فکر می‌کنم این به اتاق بستگی دارد: وقتی دستور می‌رسد چه کسی در آن است، و آیا گفتن «نه» در آنجا به قیمت جان تمام نمی‌شود. این‌قدرش را می‌توانیم بسازیم.

در نوشتن این کتاب از مدل‌های زبانی زیاد استفاده کردم، از جمله همان مدلی که درون خط لولهٔ هدف‌گیری حملات هوایی Maven در فصل~\ref{sec:1} نشسته بود؛ یادداشتی در پایان کتاب می‌گوید به نظرم این درهم‌تنیدگی با استدلال چه کرد. در ضمن، من هنوز خواندن این کتاب را تمام نکرده‌ام. مشکل این است که کلمه‌های لعنتی‌اش زیادی زیاد است، می‌فهمید که. به‌زودی تمامش می‌کنم؛ فقط به نظرم مهم‌تر بود که پیش از انتخابات میان‌دوره‌ای آمریکا منتشرش کنم. بخش عمدهٔ تاریخچهٔ بازنگری‌ها به‌رایگان در \url{https://github.com/iterabloom/antifascist.intelligence} در دسترس است.
